\documentclass[journal]{IEEEtran}
\usepackage[a5paper, margin=10mm, onecolumn]{geometry}
\usepackage{tfrupee}

\setlength{\headheight}{1cm}
\setlength{\headsep}{0mm}

\usepackage{gvv-book}
\usepackage{gvv}
\usepackage{cite}
\usepackage{amsmath,amssymb,amsfonts,amsthm}
\usepackage{algorithmic}
\usepackage{graphicx}
\usepackage{textcomp}
\usepackage{xcolor}
\usepackage{txfonts}
\usepackage{enumitem}
\usepackage{mathtools}
\usepackage{gensymb}
\usepackage{comment}
\usepackage[breaklinks=true]{hyperref}
\usepackage{tkz-euclide}

\graphicspath{{./figs/}}

\begin{document}
\title{12.88}
\author{AI25BTECH11010 - Dhanush Kumar}
\maketitle
\renewcommand{\thefigure}{\theenumi}
\renewcommand{\thetable}{\theenumi}

\textbf{Question}

Consider a discrete-time LTI system with input  
$x[n] = \delta[n] + \delta[n - 1]$  
and impulse response  
$h[n] = \delta[n] - \delta[n - 1]$.  

The output of the system will be given by:
\begin{multicols}{2}

\begin{enumerate}
    \item[(a)] $\delta[n] - \delta[n - 2]$
    \item[(b)] $\delta[n] - \delta[n - 1]$
    \item[(c)] $\delta[n - 1] + \delta[n - 2]$
    \item[(d)] $\delta[n] + \delta[n - 1] + \delta[n - 2]$
\end{enumerate}
\end{multicols}

\textbf{Solution}

For an LTI system, the output is obtained by convolution:
\begin{align}
y[n] &= x[n] * h[n].
\end{align}

Given input and impulse response:
\begin{align}
x[n] &= \delta[n] + \delta[n-1], \\
h[n] &= \delta[n] - \delta[n-1].
\end{align}

Representing them as vectors:
\begin{align}
\vec{x} &= 
\myvec{x[0] \\ x[1]}
= 
\myvec{1 \\ 1}, &
\vec{h} &= 
\myvec{h[0] \\ h[1]}
= 
\myvec{1 \\ -1}.
\end{align}

Determine output length:
\begin{align}
N_y &= N_x + N_h - 1 = 2 + 2 - 1 = 3.
\end{align}

Form the convolution (Toeplitz) matrix using $\vec{h}$:
\begin{align}
\vec{H} &= 
\myvec{
1 & 0 \\
-1 & 1 \\
0 & -1
}.
\end{align}

Compute the output vector:
\begin{align}
\vec{y} &= \vec{H}\vec{x} \nonumber \\
&= 
\myvec{
1 & 0 \\
-1 & 1 \\
0 & -1
}
\myvec{1 \\ 1} \nonumber \\
&= 
\myvec{
1(1) + 0(1) \\
-1(1) + 1(1) \\
0(1) - 1(1)
} =
\myvec{1 \\ 0 \\ -1}.
\end{align}

Thus, the output samples are:
\begin{align}
y[0] &= 1, &
y[1] &= 0, &
y[2] &= -1.
\end{align}

Hence, the output signal is:
\begin{align}
y[n] &= \delta[n] - \delta[n-2].
\end{align}

Final answer:
\begin{align}
\boxed{y[n] = \delta[n] - \delta[n-2] \quad \text{(Option (a))}}
\end{align}

\end{document}

